\documentclass[11pt,a4paper]{article}

% --- Packages ---
\usepackage[margin=.75in]{geometry}
\usepackage[utf8]{inputenc}
\usepackage[T1]{fontenc}
\usepackage{lmodern} % Fix for font not found error
\usepackage{titlesec}
\usepackage{enumitem}
\usepackage{hyperref}
\usepackage{xcolor}
\usepackage{fancyhdr}
\usepackage{comment}
\usepackage{cleveref}
\usepackage[bottom]{footmisc}
\usepackage{graphicx}
\usepackage{longtable}
\usepackage{booktabs}
\usepackage{array}
\usepackage{calc}
\usepackage{etoolbox}

% Standard Pandoc setup
\providecommand{\tightlist}{%
  \setlength{\itemsep}{0pt}\setlength{\parskip}{0pt}}

\crefformat{section}{\S#2#1#3}
\crefformat{subsection}{\S#2#1#3}
\crefformat{subsubsection}{\S#2#1#3}
\crefformat{figure}{#2Figure~#1#3}
\crefformat{footnote}{#2\footnotemark[#1]#3}

\linespread{0.9} % Riduce lo spazio tra le linee

% --- Colors ---
\definecolor{primary}{RGB}{0, 0, 0}
\definecolor{links}{HTML}{0000EE}
\definecolor{blue}{HTML}{26428B}

\hypersetup{
    colorlinks=true,
    linkcolor=links,
    urlcolor=links,
    citecolor=links,
}

% --- Section Formatting ---
\titleformat{\section}
  {\normalfont\Large\bfseries\color{blue}}{\thesection}{1em}{}
\titleformat{\subsection}
{\normalfont\large\bfseries\color{blue}}{\thesubsection}{1em}{}
\titleformat{\subsubsection}
{\normalfont\bfseries\color{blue}}{\thesubsubsection}{1em}{}

% --- Header and Footer ---
\pagestyle{fancy}
\fancyhead[L]{\small Giuseppe Capasso}
\fancyhead[R]{\small virtual-machine}
\fancyfoot[C]{\thepage}
\renewcommand{\headrulewidth}{0.4pt}

% --- Custom Title Command ---
\newcommand{\proposaltitle}[1]{
    \begin{center}
        \vspace*{0.1cm}
        {\LARGE \bfseries \color{blue} #1} \\[1.5ex]
        {\large \textbf{Giuseppe Capasso}} \\[1ex]
                \small{
            \vspace{0.1cm}
            \textbf{Materia:} virtual-machine\\
        }
            \end{center}
    \vspace{0.3cm}
    \hrule height 0.5pt
    \vspace{0.5cm}
}

\begin{document}

\proposaltitle{Setup a linux virtual machine}

\section{Setup a linux virtual
machine}\label{setup-a-linux-virtual-machine}

We will be using \href{https://www.kali.org/}{Kali Linux} for the
cybersecurity class. Any other linux environment will be fine (feel free
to reach out if you encounter any issues with software we use during
practical labs).

If you have already a linux machine skip the installation phase.

\subsection{Download ISO}\label{download-iso}

Get your Kali Linux ISO
\href{https://www.kali.org/get-kali/\#kali-installer-images}{here}.

Kali

\subsubsection{Hypervisor setup}\label{hypervisor-setup}

Now you need an hypervisor (something that runs virtual machines). We
will use Virtual Box which can be download from the
\href{https://www.virtualbox.org/wiki/Downloads}{official site} (most
likely you will have to download the Windows Hosts version.)
\includegraphics[width=5.83333in,height=2.83772in]{media/rId26.png}

Create a new virtual machine, as in pictures.
\includegraphics[width=4.16667in,height=2.77778in]{media/rId29.webp}

Give it a name and select the ISO you just downloaded.
\includegraphics[width=5.83333in,height=2.24121in]{media/rId32.png}

Finally, give it at least 2GB ram and 2CPUs and start the installation!

\subsubsection{Installation guide}\label{installation-guide}

Installation is very simple. Just follow the instructions!

First, you will see a screen like this:

Create VM

Choose graphical install and go on, you will be asked for general
informations about languages and to create a user (provide username and
password).

Create VM

Create a simple partition scheme and confirm choices (select
\texttt{yes}):

\includegraphics[width=5.83333in,height=4.375in]{media/rId42.png}
\includegraphics[width=5.83333in,height=4.375in]{media/rId45.png}

Leave blank here and click \texttt{Continue}:
\includegraphics[width=5.83333in,height=4.375in]{media/rId48.png}

Select default option for installation:

Create VM

When the installation end, you should reboot in your new virtual
machine.

\end{document}
